\section{The Inverse Trigonometric Functions}
\label{sec:inverse-trig-funct}
\markright{\sc{}The Differentials of the Inverse Trigonometric Functions}
% \baselineskip=1.5cm
% \vskip-.35in

\subsection*{The Witch of Agnesi and Inverse Tangent Function }
\begin{wrapfigure}[]{l}{1.7in}
  \captionsetup{labelformat=empty}
  \centerline{\href{https://mathshistory.st-andrews.ac.uk/Biographies/Agnesi/}{\includegraphics*[height=2in,width=1.7in]{../Figures/AgnesiPortrait}}}
  \caption{\centerline{Maria Ga\"{e}tana Agnesi}\\ \centerline{(1718-1799)}}
  \label{fig:AgnesiPortrait}
\end{wrapfigure}
In a time when women had very few options in life
\href{https://mathshistory.st-andrews.ac.uk/Biographies/Agnesi/}{Maria
  Ga\"{e}tana Agnesi} (1718-1799) was both exceptional and very
lucky. Her brilliance and talent were recognized early and nurtured by
her wealthy father, who encouraged her studies and provided her with
the best possible tutors to develop her talents. She mastered the
Calculus of Newton and Leibniz and in $1748$ wrote a series of
textbooks on the topic titled \emph{\href{https://www.maa.org/press/periodicals/convergence/mathematical-treasures-maria-agnesis-analytical-institutions}{Instituzioni Analitiche ad Uso
  Della Giovent\`u{} Italiana}\aside{This is usually translated as
    \pubtitle{Analytical Methods for the use of Italian Youth}}.} It
was immediately recognized as a masterpiece of mathematical
exposition and  was used throughout Europe during eighteenth century. It
was by far the most popular Calculus textbook in use at the time.  An
Englishman,
\href{https://mathshistory.st-andrews.ac.uk/Biographies/Colson/}{John
  Colson}, was so impressed with Agnesi's work that in $1760$ he took it
upon himself to learn Italian specifically so that he could translate
her text into English.

In her text Agnesi\aside{Her name is pronounced ``Awn-yes-ee,'' not
  ``Ag-nes-ee.'' } collected many of the known results of the time and
organized them for the instruction of students. The Latin name of one
of the curves that she used for instruction was \emph{versoria}
(``rope that turns a sail'') because of its shape. Agnesi translated
this, correctly, into  Italian as ``la versiera.''  Unfortunately
Colson mistook this for ``l'aversiera'' which means ``witch.'' As a
result this particular curve has been known ever since as the
Witch of Agnesi.


\begin{wrapfigure}[13]{r}{4.4in}
  \vskip-7mm{}
  \captionsetup{labelformat=empty}
  \includegraphics*[height=2.2in,width=4.4in]{../Figures/Witch1}
\label{fig:Witch1}
\end{wrapfigure}
As modern function notation was yet to be invented Agnesi
understood this curve geometrically as a particular set of
points. In the following diagram think of $x_0$ as moving from left
to right. For each value of $x_0$ draw the line
from the origin to the point $(x_0,1)$ and locate the vertical
coordinate $y_0$ where this line intersects the circle,
$x^2+ \left( y-\frac12 \right)^2=\left(\frac12\right)^2.$ Every
point $(x_0,y_0)$ displayed in the
diagram  is a point on the Witch of Agnesi. The Witch itself  is the
curve shown in red.

\pagebreak[4]
\begin{embeddedproblem-1line}{}
  In the diagram above show that the coordinates $x_0$ and $y_0$
  satisfy the equation, $y_0=\dfrac{1}{1+x_0^2}$. Thus the Witch is
  the graph of
  $$
  y=\frac{1}{1+x^2}.
  $$
% \hint{Notice that $(x_0, y_0)$  is not the point of intersection of the circle and the line, but they do have the
% same $y$ coordinate. It might be  helpful to label this intersection
% point $(x_1, y_0)$ and notice  that you are
% trying to find $y_0$.}
  \vskip-6mm{}
\end{embeddedproblem-1line}

%%% \centerline{  \includegraphics*[height=2.5in,width=4in]{../Figures/Witch2}}

Whenever you encounter a new curve an important question to ask is,
``What is its derivative?'' In this case you already have everything you
need to find the answer.

\begin{embeddedproblem-1line}
  Show that if $y=\frac{1}{1+x^2},$ then
  $\dfdx{y}{x}=\frac{-2x}{(1+x^2)^2}.$
\end{embeddedproblem-1line}

\begin{wrapfigure}[10]{r}{3in}
  \vskip-2mm{}
\captionsetup{labelformat=empty}
\centerline{\includegraphics*[height=1.5in,width=3in]{../Figures/Witch3}}
\label{fig:Witch3}
\end{wrapfigure}
Computing the derivative of a function as a \emph{formula} is very
useful. This has in fact been our primary focus so far in this
text. However to truly understand the relationship between a function
and its derivative nothing can replace seeing both of them graphed
together.  Look closely at the relationship between the Witch and its
derivative in the sketch at the right. Is it clear how these graphs
are related?

Notice in particular that when $x=0$ the $y$ coordinate of the Witch
is at its highest point (maximum value), whereas the $y$ coordinate of
$\dfdx{(\text{Witch})}{x}$ is zero.



% \begin{wraptable}{r}{2.4in}
%   \captionsetup{labelformat=empty}
%   \begin{tabular}{ |c|c| }
      %       \hline
      %       \multicolumn{2}{|c|}{} \\
      %       \multicolumn{2}{|c|}{\Large\bf{}Differentiation of } \\
      %       \multicolumn{2}{|c|}{} \\
      %       \multicolumn{2}{|c|}{\Large\bf{} Inverse Trig Functions} \\
      %       \multicolumn{2}{|c|}{} \\
      %       \hline
      %       $y$   \amp{}$\dx{y}$\\\hline\hline
      %       $\inverse\sin (x)$ \amp{}    $\dfrac{1}{\sqrt{1-x^2}}\dx{x} $\bigstrut{}\\
      %       $\inverse\cos (x)$ \amp{} $\dfrac{-1}{\sqrt{1-x^2}}\dx{x}$ \bigstrut{}\\
      %       $\inverse\tan (x)$\amp{}$\dfrac{1}{1+ x^2}\dx{x}$\bigstrut{}\\
      %       $\inverse\cot (x)$\amp{}$\dfrac{-1}{1+x^2}\dx{x}$\bigstrut{}\\
      %       $\inverse\sec (x)$\amp{}$\dfrac{1}{\abs{x}\sqrt{x^2-1}}$\bigstrut{}  \\
      %       $\inverse\csc x$\amp{}$\dfrac{-1}{\abs{x}\sqrt{x^2-1}}\dx{x}$\bigstrut\\\hline
      %     \end{tabular}
      %       \caption{\textcolor{blue}{\bf{} Memorize these!}}
      %       \label{tab:InvTrigDiffRules}
      %       \end{wraptable}
      % %     \end{center}
We know that the derivative of a curve at a given
point gives us the slope of the tangent line at that point. Thus at
its highest point the tangent line of the Witch will be horizontal.
That is, its derivative will be zero. This is exactly what these two
graphs are showing us.

When $x$ is just to the  right of zero, the slope of the
Witch is close to zero but negative. Similarly just to the left of zero the slope
of the Witch is close to zero but positive. At the extreme right end of the
graph above the slope of the Witch is close to zero but negative and at the
extreme left it is close to zero but positive. And all of this is reflected
in the shape of the blue (derivative) curve.

It is a little more interesting to reverse this process. That is,
suppose we could see only the derivative of the Witch (blue
curve). Could we figure out the shape of the Witch from this?

Sure we can. Well, almost.

Consider: From the extreme left of the graph to $x=0$ the
$y-$coordinate of the (blue) derivative curve is above the $x-$axis. Therefore
the slope of the Witch is positive. Clearly wherever the slope of a
function is positive the function is increasing, so we can deduce that
the Witch increases from left to right until we reach $x=0$ where the
blue derivative function crosses the $x$-axis.

Thereafter the $y-$coordinate of the blue derivative curve is below the
$x-$axis. That is, the slope of the Witch is negative, so the Witch is
decreasing.

Just by looking at the derivative curve we can see that the Witch
of Agnesi increases from left to right until we reach $x=0,$ and after
that it decreases. Unfortunately, this is not enough to completely
describe the Witch of Agnesi. While we can tell a lot about a curve by
looking at the graph of its derivative, we can't tell everything.

\begin{ProficiencyDrill}{}
  \label{drill:NonUniqueAntiderivative}
  In the sketch below the red curve is the graph of the Witch and the
  blue curve is the graph of its derivative. Convince yourself that
  the blue curve could also be the graph of the derivative of any of
  the other curves as well. As clearly as you can, explain what this
  suggests about the relationship between a
  function and its derivative.\\
  \centerline{\includegraphics*[height=1.5in,width=3in]{../Figures/Witch11}}
  
\end{ProficiencyDrill}


Next we'd like to apply this same sort of reasoning to the Witch
itself.  We'd like to answer the question, ``What curve is the Witch
of Agnesi the derivative of?''  Whatever that curve is, we'll call it
the \term{antiderivative} of the Witch for obvious reasons.

%\centerline{\includegraphics*[height=1.5in,width=3in]{../Figures/Witch6}}
\centerline{\includegraphics*[height=1.5in,width=3in]{../Figures/Witch6}}
As before, by simply looking at the graph of the Witch we can see
that: (1) At the extreme ends of its graph the antiderivative will
have a \emph{positive} slope which gets closer and closer to zero as
we go farther from the origin, but (2) since the Witch never crosses
the $x$-axis the antiderivative will never have a horizontal slope,
and (3) at $x=0$ the slope of the antiderivative will be one. A sketch
of the antiderivative of the Witch begins to emerge when we put
these three observations into our graph:\\
\centerline{\includegraphics*[height=1.5in,width=3in]{../Figures/Witch7}}

Continuing to fill between the black segments in the same  fashion
the sketch becomes even clearer:\\ 
\centerline{\includegraphics*[height=1.5in,width=3in]{../Figures/Witch8}}
Finally, connecting the dots with a smooth curve we see the the
graph of the antiderivative of the Witch of Agnesi must look like this:\\
\centerline{\includegraphics*[height=1.5in,width=3in]{../Figures/Witch9}}

We've been a little lazy with our use of language.  As we saw in
Drill~\ref{drill:NonUniqueAntiderivative}, the antiderivative is not
unique, so it is improper to speak of \emph{the} antiderivative. Any
one of the curves shown below could be \emph{an} antiderivative of the
Witch of Agnesi. (Except the black one, of course. That's the Witch.)\\
\centerline{\includegraphics*[height=1.7in,width=3in]{../Figures/arctan1}}


\subsection{The Differentials of the Inverse Trigonometric Functions}
\label{sec:diff-invers-invers}

\begin{wrapfigure}[6]{r}{2.5in}
  \vskip-3mm{}
\captionsetup{labelformat=empty}
\centerline{\includegraphics*[height=.8in,width=2.5in]{../Figures/arctan6}}
\label{fig:arctan6}
\end{wrapfigure}
These graphs should look very  familiar to you. When you learned
Trigonometry you would have encountered the inverse trig functions,
$\inverse\sin(x)$, $\inverse\cos(x)$, $\inverse\tan(x)$,
$\inverse\cot(x)$, $\inverse\sec(x)$, $\inverse\csc(x)$. In particular
the graph of the inverse tangent, seen at the right,  looks very much
like the derivative of the Witch of Agnesi, doesn't it? Do you suppose
it is possible that the derivative of the
arctangent function is the Witch of Agnesi\aside{Of course it is. Why
  else would we have led you down this path? But there are some
  subtleties here that we shouldn't ignore. We will proceed cautiously}?

  Reasoning from pictures, as we've just done, can be very useful. It
  gives us strong insights into problems and we encourage you to do it
  whenever you can.

  But.  The most we can hope for when reasoning from pictures is an
  intuitive feel for the problem. We always need to confirm our
  intuition analytically. \emph{Always}.

So we will verify that our conjecture -- that the derivative of the
arctangent function is the Witch of Agnesi -- analytically.

The function name \emph{arctangent} literally means ``the angle (arc)
whose tangent is\aside{It can be helpful to think of the arctangent
  like this: When you see $y=\arctan(x)$ think (or say) ``$y$ is the
  angle whose tangent is $x$.''  }.''  This means that $ y=\arctan(x)$
precisely when $\tan(y)=x$.  Differentiating $\tan(y)=x$ gives:
\begin{align*}
  \dx(\tan(y)) \amp{}= \dx{x}\\
  \sec^2(y)\dx{y} \amp{}= \dx{x}\\
  \dx{y}\amp{}= \frac{1}{\sec^2(y)}\dx{x}.
\end{align*}
This is a correct formula for the differential of the arctangent
function. However, there are two problems: (1) it does not seem to
match our conjecture, and (2) even if it is correct this formula will
not be very useful to us in this form since we have $\dx{y}$
in terms of $y$ itself. We can address both of these problems by
finding $\dx{y}$ in terms of $x$ and $\dx{x}$.

\begin{wrapfigure}[]{r}{1.7in}
  \vskip-6mm{}
  \captionsetup{labelformat=empty}
  \centerline{\includegraphics*[height=1.7in,width=1.8in]{../Figures/arctanMnemonic2}}
  \label{fig:arctanMnemonic2}
\end{wrapfigure}
It the diagram at the right we see that $\tan(y) = \frac{x}{1}=x$, so
that $y$ is ``the angle whose tangent is $x$'' or $y=\arctan(x)$. Also
recall  that  the hypotenuse of $\triangle
OAB$ is the ``segment that cuts the circle,''
or  $\sec(y)$. From the Pythagorean Theorem we see that
$\sec^2(y)=1+x^2$ 
% is ``the
% angle whose tangent is $x$.''  Since the secant function is equal to
% $\dfrac{\text{hypotenuse}}{\text{adjacent}}$, the Pythagorean Theorem
% tells us that
% $$
% \sec(y)= \frac{\sqrt{1+x^2}}{1}=\sqrt{1+x^2},
% $$
so that
$$
\dx{y}=\frac{1}{\sec^2(y)}\dx{x}=\frac{1}{1+x^2}\dx{x},
$$
or
$$
\dfdx{y}{x}=\frac{1}{1+x^2}
$$
So our intuition was correct. The Witch is indeed the derivative of
the arctangent function.  A similar derivation can be used for the
arccotangent.

%\pagebreak[4]

\begin{embeddedproblem-enumerate}{}
  \begin{enumerate}[label={\bf{}(\alph*)}]
  \item the arccotangent function is defined as  $y=\arccot(x)$ if and
    only if $\cot(y)=x$. Proceed as we did above, to show that
    $$
    \dx{\left(\arccot(x)\right)} = \frac{-1}{1+x^2}\dx{x}.
    $$
  \item Derive the same result from  the identity
    $\arccot(x)=\pi/2-\arctan(x)$.
  \end{enumerate}
\end{embeddedproblem-enumerate}

\pagebreak[4]
\begin{embeddedproblem-enumerate}
  \begin{enumerate}[label={\bf{}(\alph*)}]
  \item Show that if $y=\arctan(x)$ then $y$ satisfies the second order
    differential equation
    $$
    (1+x^2)\dfdxn{y}{x}{2}+2x\dfdx{y}{x}=0.
    $$
  \item The function   $y=\arccot(x)$ satisfies the same differential
    equation. Show this in two different ways.
    \begin{enumerate}[label={\bf{}(\roman*)}]{}
    \item By direct computation, just as you did part (a).
    \item By direct computation, after first observing that $\arccot(x)=\frac\pi2-\arctan(x)$.
    \end{enumerate}
  \end{enumerate}{}
\end{embeddedproblem-enumerate}


The arcsine and arccosine functions are defined in a similar
manner. If $y=\arcsin(x)$ or $y=\arccos(x)$ then $y$ is ``the angle
whose sine (cosine) is $x$.''
To find the  differentials of $y=\arcsin(x)$ and $y=\arccos(x)$ we
proceed just as we did with the $\arctan(x)$ and $\arccot(x)$:
\begin{align*}
  \dx{(\sin(y))}\amp{}= \dx{x}\\
  \cos(y)\amp{}=\dx{x}\\
  \dx{y}\amp{}=\frac{1}{\cos(y)}\dx{x}
\end{align*}

\begin{ProficiencyDrill}
  \label{drill:InvSin1}
  Use a  diagram like the one above  below to show that
  $\dx(\arcsin(x))=\frac{\pm 1}{\sqrt{1-x^2}}$.
  %\\
%    \centerline{\includegraphics*[height=1.5in,width=2.4in]{../Figures/CosInvSin1}}
\end{ProficiencyDrill}

\begin{embeddedproblem-enumerate}{}
  \begin{enumerate}[label={\bf{}(\alph*)}]
  \item Use the fact that $y=\arccos(x)$ if and only if $\cos(y)=x$
    and proceed as we did above to show that
    $$
    \dx\left(\arccos(x)\right)= \frac{\mp 1}{\sqrt{1-x^2}}\dx{x}.
    $$
    \comment{Notice how this formula differs from the one you found
      in Drill~\ref{drill:InvSin1}.}
  \item Derive the same result from the identity:  $\arccos(x)=
    \frac{\pi}{2}-\arcsin(x)$.
  \end{enumerate}
\end{embeddedproblem-enumerate}


      %       \pagebreak[4]
\begin{embeddedproblem-enumerate}
  \begin{enumerate}[label={\bf{}(\alph*)}]
  \item Show that if $y=\arcsin(x)$ then $y$ satisfies the second order
    differential equation
    $$
    (1-x^2)\dfdxn{y}{x}{2}-x\dfdx{y}{x}=0.
    $$
  \item The function   $y=\arccos(x)$ satisfies the same differential
    equation. Show this in two different ways.
    \begin{enumerate}[label={\bf{}(\roman*)}]{}
    \item By direct computation, just as you did part (a).
    \item By direct computation, after first observing that $\arccos(x)=\frac\pi2-\arcsin(x)$.
    \end{enumerate}
  \end{enumerate}{}
\end{embeddedproblem-enumerate}


\begin{wrapfigure}[]{r}{2in}
  \captionsetup{labelformat=empty}
  \centerline{\includegraphics*[height=1.5in,width=2in]{../Figures/InvSecTriangle}}
  \label{fig:InvSecTriangle}
\end{wrapfigure}

To find the differential of $\arcsec(x)$ we refer to the sketch at the
right and we proceed exactly as before. If $y=\arcsec(x)$, then
$x=\sec(y)$ so
\begin{align*}
  \dx{(\sec(y))}      \amp{}=\dx{x}\\
  \sec(y)\tan(y)\dx{y}\amp{}= \dx{x}
                        \intertext{so that}
                        \dx{y}              \amp{}=\frac{{1}}{\sec(y)\tan(y)}\dx{x}.
\end{align*} 

\begin{ProficiencyDrill-1line} Show that: $ \dx\left(\arcsec(x)\right)=\frac{\pm1}{x\sqrt{x^2-1}}.  $
\end{ProficiencyDrill-1line}

\begin{embeddedproblem-enumerate}{}
  \begin{enumerate}[label={\bf{}(\alph*)}]
  \item Use the fact that $y=\arccsc(x)$ if and only if $\csc(y)=x$
    and proceeding as we did above, show that
    $$
    \dx\left(\arccsc(x)\right)= \frac{\mp 1}{\sqrt{1-x^2}}\dx{x}.
    $$
  \item Derive the same result from the identity:  $\arccsc(x)=
    \frac{\pi}{2}-\arcsec(x)$.
  \end{enumerate}
\end{embeddedproblem-enumerate}


\begin{embeddedproblem-enumerate}
  \begin{enumerate}[label={\bf{}(\alph*)}]
  \item Show that if $y=\arcsec(x)$ then $y$ satisfies the second order
    differential equation
    $$
    x(x^2-1)\dfdxn{y}{x}{2}+(2x^2-1)\dfdx{y}{x}=0.
    $$
  \item The function   $y=\arccsc(x)$ satisfies the same differential
    equation. Show this in two different ways.
    \begin{enumerate}[label={\bf{}(\roman*)}]{}
    \item By direct computation, just as you did part (a).
    \item By direct computation, after first observing that $\arccsc(x)=\frac\pi2-\arcsec(x)$.
    \end{enumerate}
  \end{enumerate}{}
\end{embeddedproblem-enumerate}


We now have the differentials of all of the inverse trigonometric
functions as seen in the table below.
\begin{center}
  \begin{tabular}{|cc|} \hline
    Function\amp{}Derivative\\\hline\hline
      $y=\arctan(x)$\amp{}$\displaystyle\dfdx{y}{x}=\frac{1}{\sec^2(y)}=\frac{1}{1+x^2}$\\
      \amp{}\\
      $y=\arccot(x)$\amp{}$\displaystyle\dfdx{y}{x}=\frac{-1}{\csc^2(y)}=\frac{-1}{1+x^2}$\\
      \amp{}\\
      $y=\arcsin(x)$\amp{}$\displaystyle\dfdx{y}{x}=\frac{1}{\cos(y)}=\frac{\pm1}{\sqrt{1-x^2}}$\\
      \amp{}\\
      $y=\arccos(x)$\amp{}$\displaystyle\dfdx{y}{x}=\frac{-1}{\sin(y)}=\frac{\mp1}{\sqrt{1-x^2}}$\\
      \amp{}\\
  $y=\arcsec(x)$ \amp{}$\displaystyle\dfdx{y}{x}=\frac{1}{\sec(y)\tan(y)}=\frac{\pm1}{x\sqrt{x^2-1}}$\\
      \amp{}\\
$y=\arccsc(x)$ \amp{} $\displaystyle\dfdx{y}{x}=\frac{-1}{\csc(y)\cot(y)}=\frac{\mp1}{x\sqrt{x^2-1}}$\\
      \amp{}\\\hline
    \end{tabular}
  \end{center}

% The need for the domain
% restrictions in the center column should be clear if you think about
% the nature of the trigonometric functions. For example, since
% the sine (and cosine) function only return numbers between $-1$ and
% $1$ it makes no sense to write, for example, $\arcsin(2)$ because this
% asks for ``the arc (angle) whose sine is $2$,'' and there is no such
% arc\aside{At least not among the real numbers. It is possible to find
%   a complex ``angle'' whose sine is $2$. But this is well beyond the
%   scope of this text.}.  Similarly the inverse secant and inverse
% cosecant functions are only defined for values of $x\le -1$ and
% $x\ge1$ because the secant and cosecant functions only return values
% in that range.

% % \begin{wraptable}[21]{r}{4.3in}
% %   \vskip-1mm{}
% %   \captionsetup{labelformat=empty}
% \begin{center}
%     \begin{tabular}{|ccc|} \hline Function\amp{}
%       Domain\amp{}Derivative\\\hline\hline $y=\arctan(x)$ \amp{} Any
%       real
%       number\amp{}$\displaystyle\dfdx{y}{x}=\frac{1}{\sec^2(y)}=\frac{1}{1+x^2}$\\
%       \amp{}\amp{}\\
%       $y=\arccot(x)$ \amp{} Any real number\amp{}$\displaystyle\dfdx{y}{x}=\frac{-1}{\csc^2(y)}=\frac{-1}{1+x^2}$\\
%       \amp{}\amp{}\\
%       $y=\arcsin(x)$ \amp{} $-1\le x\le 1$\amp{}$\displaystyle\dfdx{y}{x}=\frac{1}{\cos(y)}=\frac{\pm1}{\sqrt{1-x^2}}$\\
%       \amp{}\amp{}\\
%       $y=\arccos(x)$ \amp{} $-1\le x\le 1$\amp{}$\displaystyle\dfdx{y}{x}=\frac{-1}{\sin(y)}=\frac{\mp1}{\sqrt{1-x^2}}$\\
%       \amp{}\amp{}\\
%       \multirow{2}{*}{$y=\arcsec(x)$} \amp{} $x\le-1$\amp{}\multirow{2}{*}{$\displaystyle\dfdx{y}{x}=\frac{1}{\sec(y)\tan(y)}=\frac{\pm1}{x\sqrt{x^2-1}}$}\\
%       \amp{}or $x\ge1$\amp{}\\
%       \amp{}\amp{}\\
%       \multirow{2}{*}{$y=\arccsc(x)$} \amp{} $x\le-1$\amp{}\multirow{2}{*}{$\displaystyle\dfdx{y}{x}=\frac{-1}{\csc(y)\cot(y)}=\frac{\mp1}{x\sqrt{x^2-1}}$}\\
%       \amp{}or $x\ge1$\amp{}\\
%       \amp{}\amp{}\\\hline
%     \end{tabular}
%   \end{center}

% \caption{\textcolor{blue}{Differentials of Inverse Trigonometric functions}}
%    \label{tab:InvTrigDiff1}
% \end{wraptable}

\subsection{ Inverse Functions}
\label{sec:InvFunctions}


But\footnote{Held for later insertion: The need for the domain
restrictions in the center column should be clear if you think about
the nature of the trigonometric functions. For example, since
the sine (and cosine) function only return numbers between $-1$ and
$1$ it makes no sense to write, for example, $\arcsin(2)$ because this
asks for ``the arc (angle) whose sine is $2$,'' and there is no such
arc\aside{At least not among the real numbers. It is possible to find
  a complex ``angle'' whose sine is $2$. But this is well beyond the
  scope of this text.}.  Similarly the inverse secant and inverse
cosecant functions are only defined for values of $x\le -1$ and
$x\ge1$ because the secant and cosecant functions only return values
in that range.
} you have surely noticed that the information given in the
table above presents some difficulties to
computation. The presence of the ``$\pm$'' symbol suggests that when
we compute the derivative of the arcsine, arccosine, arcsecant, and
arccosecant functions we get to choose whether to choose positive or
negative value.

But this seems nonsensical, doesn't it? After all the slopes of each
of these curves at a given value of $x$ will be positive or negative
independent of any choices we might make. It is equally odd that
this plus-or-minus conundrum doesn't seem to be a problem with the
arctangent or the arccotangent functions. Why would that be?

The problem here -- or at least part of the problem -- is that we
haven't actually defined \emph{clearly} what we mean by
``$\arctan$,'' or any of the other inverse  functions. In fact we
haven't really defined them at all.
We will
address this next.

We'll begin with the arctangent because it is the easiest inverse trig
function to think about. We said earlier that $y=\inverse\tan(x)$
precisely when $x=\tan(y)$, but what does this really mean?

To make this as concrete as possible observe that
$\tan\left(\frac{\pi}{4}\right)=1$ and therefore
$\arctan(1)=\frac{\pi}{4}$. This means that the numbers $1$, and
$\frac{\pi}{4}$ are associated by the tangent and arctangent
functions. If $1$ goes into the tangent function, $\frac{\pi}{4}$
comes out. If $\frac{\pi}{4}$ goes into the arctangent function, $1$
comes out.  This means that if we were to graph the tangent function
one point on the graph would be $\left(\frac{\pi}{4}, 1\right)$. If we
were to graph the arctangent function one point on the graph would be
$\left(1, \frac{\pi}{4}\right)$. Although they are different functions
the tangent and arctangent function contain the same information:
$1$ and $\frac{\pi}{4}$ are associated.

Moreover, to switch from one function to the other all we have to do
is swap the roles of $1$ and $\frac{\pi}{4}$.

More generally, if $x$ and $y$ are associated by the tangent and
arctangent functions (that is, if $y=\inverse\tan(x)$ then
$x=\tan(y)$) then to switch from one function to the other all we have
to do is swap the roles of $x$ and $y$. The set of all things that
``go into'' a function is called its \term{domain}, and the set of all
things that ``come out'' of a function is called its \term{range}.

So, at last, we see that we can switch from tangent to arctangent by
swapping the roles of the domain and range. That is, the domain of
$\inverse\tan(x)$ is the range  of $\tan(x)$ and the range of
$\inverse\tan(x)$ is the domain of $\tan(x)$.

We graph a function by plotting all of the ordered pairs of numbers
which are associated by the function. Usually we take the horizontal
axis to be the domain and the vertical axis to be the range. Thus, if
we have the 


%\centerline{\includegraphics*[height=2in,width=5in]{../Figures/arctan3}}
\centerline{\includegraphics*[height=1.5in,width=4in]{../Figures/arctan3}}
The sketch above shows the graph of $y=\tan(x)$ (in red) and the same
graph (in blue) reflected about the graph of the line $y=x$ (dashed, in
black). This \emph{ought} to be the graph of $y=\arctan(x)$.  But it
takes only a few seconds consideration to realize that the blue sketch
on the right can't possibly be the graph of \emph{any}
function. Consider, for example $\arctan(0)$. Is this equal to $0$? Or
is it $\pi$? $2\pi$?  Or maybe $\arctan(0)=-\pi$?


That we seem to have a choice clearly violates what we mean we say
``function'' since a mathematical function returns \emph{exactly} one
output for a given input. There is no choice. So the sketch on the
right can't possibly be the graph of a \term{function}. It must be
something else. But what?

The sketch on the right is the graph of the \term{multifunction}
$\arctan(x)$. Each of the blue curves on the right is one
\term{branch} of the multifunction, $\arctan(x)$. In fact all of the
``arc'' functions from trigonometry are multifunctions.
Mult-functions are interesting objects and are well worth studying.
But this is not the time for that study. Right now we are only
interested in the properties of the \term{inverse function} of
$\tan(x)$. Since $\arctan(x)$ is not a function, it can't be the
inverse function of $\tan(x)$ so we will stop thinking about
it now.

But if $\arctan(x)$ is not the inverse of $\tan(x)$ what is?  Could it
be as simple as choosing just one of the branches of the
$\arctan(x)$, say the one between $y=-\frac{\pi}{2}$ and $y=\frac{\pi}{2}$ and calling it the the inverse of  $\tan(x)$?\\
\centerline{\includegraphics*[height=2.7in,width=4.5in]{../Figures/arctan5}}
Actually, yes. It is just that simple. And the usual choice is the one
we've indicated.  In this text we will designate this branch of
$\arctan(x)$ with ``$\inverse\tan(x)$.'' Thus ``$\inverse\tan(x)$''
represents the \term{inverse tangent} function, while $\arctan(x)$
represents a particular multifunction that we will henceforth ignore
as much as possible.

But what is $\inverse\tan(x)$ the inverse of? From the sketch above it
should be clear that restricting the \alert{range} of
$\inverse\tan(x)$ to values in the interval $(-\pi/2, \pi/2)$ forced
us to restrict the \alert{domain} of $\tan(x)$. So
$$
\inverse\tan(x) \text{ with domain all real numbers, } (\RR)
$$
is the inverse of
$$
\tan(x) \text{ with domain, } (-\pi/2, \pi/2),
$$
and it is \emph{not} the inverse of
$$
\tan(x) \text{ with domain, } \RR
$$
because that function  has no inverse\aside{This is because $$\tan(x) \text{ with domain } \RR$$
  is \alert{not} same function as
  $$\tan(x) \text{ with domain } -\frac{\pi}{2}\lt{}x\lt{}\frac{\pi}{2}.$$  A
  function has two parts: (1) The set of inputs
  (domain) and (2) the rule associating input with output. Most of the
  time the domain is the set of all real numbers ($\RR$), so we don't
  bother to explicitly state that the domain of, say $f(x)=x^2$ is
  $\RR${}. Unless otherwise
  specified we simply  assume that it is.\\
  However, if we change either part of a function then we've created a
  new function so $f(x)=x^2$ is \alert{not} the same function as
  $f(x)=x^2$ with domain $0\le x\le 1$.  Properly speaking, we should \emph{always}
  specify the domain of our functions, but we don't. The custom is to
  assume that the domain of a function is $\RR$ unless otherwise
  stated.\\
  But we (mathematicians) add to the confusion when dealing with
  inverse functions because we habitually violate this custom as well.\\
  The function, $\tan(x)$ cannot be inverted because it has more than
  one branch. Because we want and need an inverse tangent function we
  restrict the range of $\inverse\tan(x)$. But now $\inverse\tan(x)$
  is \alert{not} the inverse of $\tan(x)$, which has too many
  branches. It is the inverse of
  $\tan(x), \text{ with domain } -\frac{\pi}{2}\lt{}x\lt{}\frac{\pi}{2}$ which
  has only one
  branch. \\
  To use our notation as it was intended to be used we \emph{should}
  say that $\inverse\tan(x)$ is the inverse of $\tan(x)$, with domain
  $-\frac{\pi}{2}\lt{}x\lt{}\frac{\pi}{2}$, which is true. But instead, we
  typically just say that it is the inverse of $\tan(x)$, which
  strictly speaking, is not true. \\
  When we (mathematicians) talk amongst ourselves this is not a
  problem. We all understand what we mean. But by definition, students
  usually have only a tenuous grasp of these abstract distinctions so
  it is a bit unfair of us (mathematicians) to speak to students in
  what amount to incomplete sentences. \\
  We (the authors) apologize.}.


To be clear we could have chosen any of the branches of $\arctan(x)$
to be $\inverse\tan(x)$. The choice is (almost) completely
arbitrary. We've made this particular choice because (1) it seems to
us to be the most natural, (2) it has been the standard choice for a long
time, and (3) most existing software has this choice built in.

\begin{ProficiencyDrill}
  Suppose we had chosen
  $\frac{\pi}{2}\lt{}\inverse\tan(x)\lt{}\frac{3\pi}{2}$. What function (with
  domain) is that the inverse of?
\end{ProficiencyDrill}

      %       The figure above shows the graph of $\tan(x)$ and its inverse,
      %       $\inverse\tan(x)$. But notice that we've explicitly constrained
      %       the domain of
      %       $\tan(x)$. Do you see why we had to do that?  We can only include one
      %       branch of $\tan(x)$ because $\inverse\tan(x)$ has only one branch of
      %       $\arctan(x)$ and they have to match up if they are going to be
      %       mutually inverse.


      %       Each of the curves in this graph is called a branch of the arctangent.
      %       The graph in red above is called the \term{principle
      %       branch} because, by convention, it is the one chosen by default --
      %       unless there is a compelling reason to choose another.  This is the one where
      %       $-\frac{\pi}{2}\lt{}y\lt{}\frac{\pi}{2}$.one.

All of this fussiness is really just about making our abstract
definitions useful  and
consistent. It would be nice if these details only impinged on us in an
abstract setting but unfortunately some  practical difficulties do
come up, as the following drill shows.

%\pagebreak[4]
% \begin{ProficiencyDrill}
%   Compare the graphs  of each of the following:
%   \begin{enumerate}[label={\bf{}(\alph*)}]
%     \begin{multicols}{4}
%     \item $y=\inverse\tan\left(\frac1x\right)$
%     \item $y=\inverse\tan\left(\frac{1}{x-1}\right)$
%     \item $y=\inverse\tan\left(\frac{1}{x-1}\right)$
%     \item $y=\inverse\tan\left(\frac{x+1}{1-x}\right)$
%     \end{multicols}
%   \end{enumerate}
%   % $y=\inverse\tan\left(\frac{1+x}{1-x}\right)$ to the graph of
%   % $y=\inverse\tan(x)+\inverse\tan(1)=\inverse\tan(x)+\frac{\pi}{4}$.\\
%   % \centerline{\includegraphics*[height=1in,width=4in]{../Figures/arctanabs1}}
%   Can you explain why they are discontinuous even though $y=\inverse\tan(x)$
%   is continuous?
% \end{ProficiencyDrill}

\begin{embeddedproblem}
  \vskip-2mm{}
  Find \emph{all} solutions of  $\tan(x)=1$ and
  $x=\inverse\tan\left(1\right)$. Do they have the same set of
  solutions?\\
  \hint{Obviously, they do not. Otherwise we wouldn't have asked the
    question. What is the difference between the two sets of
    solutions?}
\end{embeddedproblem}

\vskip-2mm{}
We obtained our branch for the single-valued function $y=\inverse\tan(x)$
from the multifunction $y=\arctan(x)$ by restricting its  range to
$-\frac{\pi}{2}\lt y\lt\frac{\pi}{2}$. A similar restriction would
allow us to obtain the single-valued
function $\inverse\cot(x)$ from the multifunction $\arccot(x)$, but it
is a bit simpler to once again use a trigonometric  identity.

\begin{ProficiencyDrill-1line}
  Use the identity $\inverse\cot(x) = \frac{\pi}{2}-\inverse\tan(x)$
  to define the domain of $y=\inverse\cot(x)$.
\end{ProficiencyDrill-1line}

\begin{embeddedproblem}
  Show that each of the following statements is true.
  \begin{enumerate}[label={\bf{}(\alph*)}]
    \begin{multicols}{2}
    \item $\dx{\left(\inverse\tan\left(\dfrac{1+x}{1-x}\right)\right)} = \dfrac{1}{1+x^2}\dx{x}$
    \item $\dx{\left(\inverse\tan\left(\dfrac{2+x}{1-2x}\right)\right)} = \dfrac{1}{1+x^2}\dx{x}$
    \item $\dx{\left(\inverse\tan\left(\dfrac{10+x}{1-10x}\right)\right)} = \dfrac{1}{1+x^2}\dx{x}$
    \item $\dx{\left(\inverse\tan\left(\dfrac{2543+x}{1-2543x}\right)\right)}
      = \dfrac{1}{1+x^2}\dx{x}$
    \end{multicols}
  \item Now show that
    $\dx{\left(\inverse\tan\left(\dfrac{\kappa+x}{1-\kappa x}\right)\right)}
    = \dfrac{1}{1+x^2}\dx{x}$ where $\kappa$ is any constant.
  \item This is weird isn't it? None of the functions on the left is
    $\inverse\tan(x)$  but they all have the same derivative as $\inverse\tan(x)$. Can you
    explain this?
    \hint{Let $y_1=\inverse\tan(x)$ and $y_2=\inverse\tan(\kappa)$. What is
      $\tan(y_1+y_2)$?}
  \end{enumerate}
  \vskip-6mm{}
\end{embeddedproblem}



\begin{wrapfigure}[12]{r}{1.3in}
  \vskip-5mm{}
  \captionsetup{labelformat=empty}
\centerline{\includegraphics*[height=1.9in,width=1.3in]{../Figures/arcsin5}}
\label{fig:arcsin5}
\end{wrapfigure}
The single valued branches for $y=\inverse\tan(x)$ and
$y=\inverse\cot(x)$, were actually fairly apparent from the graphs of
$x=\tan(y)$ and $x=\cot(y)$, despite how much time we spent talking
about them. All we had to do was choose one continuous branch from
either graph. We described the process so carefully because for the
other inverse trigonometric functions things are a little more
problematic. For those we will have to use the derivatives of our
``arc'' multifunctions as a guide, rather than the multifunctions
themselves.


When we graph $\arcsin(x)$ by flipping the graph of $\sin(x)$ across
the line $y=x$ we get the sketch at the right. Obviously $\arcsin(x)$
is also a multifunction, not a function.  For the value of $x$ shown
in the sketch the arrows point to a few of the possible values of
$\arcsin(x)$. There are again infinitely many.

Since $\arcsin(x)$ is not a function at all it can't be the inverse
\emph{function} of $\sin(x)$.  Once again, to \emph{create} an inverse
sine function (which we will call $\inverse\sin(x)$) we will have to
restrict the range of $\inverse\sin(x)$ just like we restricted the
range of $\inverse\tan(x)$.  The question is, what restriction should
we impose? When we were looking to restrict the range of the
$\inverse\tan(x)$ there was a ``natural'' choice that was clearly
visible in the graph.  That doesn't seem to be true this time. We'll
have to find some other criterion to base our choice on.

\begin{wrapfigure}[13]{l}{1.3in}
  \vskip-4mm{}
  \captionsetup{labelformat=empty}
\centerline{\includegraphics*[height=1.9in,width=1.3in]{../Figures/arcsin6}}
\label{fig:arcsin6}
\end{wrapfigure}
Looking back at our differentiation table we see that the
derivative of $\arctan(x)$ (and thus the derivative of
$\inverse\tan(x)$) is always positive while the derivative of the
$\arccot(x)$ (and thus the derivative of $\inverse\cot(x)$) is always
negative. This happened naturally with the tangent and cotangent, but
this is not true of the other trigonometric functions.

In particular, it is not true of the sine function and its inverse.
As we can see in the sketch at the left the slope of $\arcsin(x)$
is sometimes positive and sometimes negative.
Since we must restrict the range of $\inverse\sin(x)$, we may as well
make a convenient choice. From the sketch we see that if we restrict
the range of $\inverse\sin(x)$ to the interval $[-\pi/2, \pi/2]$ the
graph of $\inverse\sin(x)$ (in blue) will \emph{always}\aside{Except
  at the endpoints of course.} have a
positive slope, and therefore a positive derivative.

With this restriction in place we \emph{define} the derivative of the
inverse sine as follows.

%\pagebreak[4]
\begin{mydefinition}[The Derivative of the Inverse Sine]
  \label{def:InvSin}
 The derivative of the inverse sine function is defined to be,\\
 \centerline{$\displaystyle\dfdx{(\inverse\sin(x))}{x}=\frac{1}{\sqrt{1-x^2}}$,} where $-1\lt
 x\lt 1$.
  \vskip-4mm{}
\end{mydefinition}
With definition~\ref{def:InvSin} in place
we won't need to concern ourselves with the ``plus-or-minus'' issue
for the derivative of  $\inverse\sin(x)$. We will similarly simply
things for ourselves regarding the derivative of $\inverse\cos(x)$.

\begin{Digression}[Undefined Derivatives]
  Notice that restricting the \emph{\underline{range}} of the inverse
  sine forces us to restrict the \emph{\underline{domain}} of the sine
  itself. Because they are mutually inverse the domain of one is the
  range of the other. So $\inverse\sin(x)$ is the inverse of $\sin(x)$
  with domain $[-\pi/2, \pi/2]$, but \alert{not} the inverse of
  $\sin(x)$ with domain $\RR$, because that function does not have an
  inverse.

  It is the nature of our topic which forces us to be so very careful
  in making our definitions, not any perverse desire for complexity
  for its own sake. Sometimes things are just complex and we have to
  deal with that.  But since we are having to pay such close attention
  to the details here we may as well point out one other small
  difficulty.
  
  If we don't think about it very hard it seems like the derivative of
  a function should be defined on the same domain as the function
  itself. Every point where a function is defined it should have a
  slope, right?

  But notice that this is not true of $\inverse\sin(x)$. The domain of
  $\inverse\sin(x)$, is $[-1,1]$ but the derivative of the inverse
  sine is clearly not defined at the endpoints. (Do you see why not?)
  So the domain of the inverse sine is slightly larger than the domain
  of its derivative.  Indeed, it is often the case that the domain of
  a function is (usually slightly) larger than the domain of its
  derivative.

  This probably all seems like a lot of fussy, nit-picking on our
  part. After all, the \emph{tangent line} of $\inverse\sin(x)$ is
  perfectly well defined at the endpoints. In both cases it is a
  vertical line, so what's all the fuss about?

  It is this.  The tangent line is not the derivative. We spend a good
  deal of time at the beginning of a Calculus course talking about and
  drawing tangent lines simply because it is helpful, and intuitively
  appealing, to do so. As a result of this repetition it is easy to
  begin to identify the tangent line of the graph of a function with
  the derivative of the function. But they are two different entities.

  They must be. Recall that we used the derivative concept to define
  the tangent line in definition~\ref{def:TangentLine}. So the
  existence of the line tangent to the graph of $y=\inverse\sin(x)$ at
  $x=\pm1$ doesn't necessarily mean that the derivative is defined
  there. And in fact, in this case it is not.

  Under what conditions a derivative is, and is not, defined is not a
  matter we have taken up yet because we do not yet have all of the
  tools we'll need. We will return to this question in
  Section~\ref{subsec:underl-deriv}.
\end{Digression}

Inverting the cosine function has all of the same issues.  By flipping
the graph of the cosine across the line $y=x$ we can see the graph of
the multifunction, $\arccos(x)$. The cosine with domain $\RR$ is not
invertible so inverting the cosine consists of selecting an
appropriate branch of $\arccos(x)$ by constraining the the range of
$\arccos(x)$ (equivalently, constraining the domain of $\cos(x)$).

\begin{ProficiencyDrill}
  \begin{wrapfigure}[]{r}{2in}
    \vskip-5mm{}
\captionsetup{labelformat=empty}
  \centerline{\includegraphics*[height=1.4in,width=1.6in]{../Figures/arccos2}}
\label{fig:arccos2}
\end{wrapfigure}
  We used our freedom to constrain the range of $\inverse\sin(x)$ to
  guarantee that $\dfdx{(\inverse\sin(x))}{x}$ is always
  positive. Similarly we constraint the range of $\inverse\cos(x)$ to
  guarantee that $\dfdx{(\inverse\cos(x))}{x}$ is always negative.
  Use the sketch at the right to show that if the
  range is constrained to the interval $[0,\pi]$ then
  $$
  \dfdx{(\inverse\cos(x))}{x} = \frac{-1}{\sqrt{1-x^2}}.
  $$
\end{ProficiencyDrill}

      %       The derivative of inverse \emph{tangent} is always positive no matter which
      %       branch we choose. Similarly The derivative of
      %       inverse \alert{co}\emph{tangent} is always negative. There is no
      %       choice where the inverses of the tangent and cotangent are involved.

      %       But when we were choosing branches for $\inverse\sin(x)$ and
      %       $\inverse\cos(x)$ we found that we did have a choice. By choosing our
      %       branch carefully we can guarantee that either one has a positive or a
      %       negative derivative.  To be consistent with the inverse tangent we
      %       deliberately contrived to chose branches that made the derivative of
      %       the \emph{sine} positive and made the derivative of the
      %       \alert{co}\emph{sine} negative.

Forcing the derivative of the inverse secant to be positive is a bit
harder for a couple of reasons. The first is just that these functions
are used less and thus we are not as familiar with them.  The other is
the nature of the derivative formula itself. Recall from the
table above that
$$
\dfdx{(\arcsec(x))}{x}=\frac{\pm1}{x\sqrt{x^2-1}}.
$$
It is the presence of the ``$x$'' in the denominator \emph{along with} the
``$\pm1$'' in the numerator that is troublesome. By choosing the range
constraint judiciously we can control
whether the ``plus'' or the ``minus'' is chosen in the numerator but
no matter how we constrain the \emph{range} of $\inverse\sec(x)$ the
\emph{domain} will always include both positive and negative numbers,
so the ``$x$'' in the denominator will make the derivative of
$\inverse\sec(x)$ positive for some values and negative for others. We
don't want that. We want it 
to be positive.

Usually when we want things in the world to be simple, Nature just
laughs at us. But not this time. This time the problem is simply that
we have too many choices.  Both the ``$\pm1$'' in the numerator and
the ``$x$'' in the denominator will influence the sign (positive or
negative) of the expression $\frac{\pm1}{x\sqrt{x^2-1}}$, but if they
are both in play it is hard to control the sign. So we just declare,
by fiat, that the ``$x$'' in the denominator will be positive to
simplify things.  Specifically, we'll replace the ``$x$'' in the
denominator with $\abs{x}$. Now we can control the sign of the
derivative of the $\inverse\sec(x)$ by simply constraining its range.

\begin{embeddedproblem-enumerate}
%   \begin{wrapfigure}[]{r}{2in}
%     \vskip-10mm{}
%     \captionsetup{labelformat=empty}
%     \centerline{\includegraphics*[height=1.6in,width=2in]{../Figures/arcsec5}}
%     \label{fig:arcsec5}
%   \end{wrapfigure}
  \begin{enumerate}[label={\bf{}(\alph*)}]
  \item   Use the sketch below \\
        \centerline{\includegraphics*[height=1.6in,width=2in]{../Figures/arcsec5}}
        to show that if we constrain the range of the
    $\inverse\sec(x)$ to be the interval $(0, \pi/2)\cup(\pi/2,\pi)$ then 
    $$
    \dfdx{(\inverse\sec(x))}{x}=\frac{1}{\abs{x}\sqrt{x^2-1}}.
    $$
  \item Use the identity $\inverse\csc(x)=\frac\pi2-\inverse\sec(x)$ to
    show that
    $$
    \dfdx{(\inverse\csc(x))}{x}=\frac{-1}{\abs{x}\sqrt{x^2-1}}.
    $$
  \item What is the range of $\inverse\sec(x)$?
  \end{enumerate}
  \vskip-3mm{}
\end{embeddedproblem-enumerate}


% \begin{embeddedproblem-enumerate}
%   \begin{enumerate}[label={\bf{}(\alph*)}]
%   \item Sketch the graph of the multifunction $\arcsec(x)$ and use
%     it to identify the range constraint that guarantees that
%     $\dfdx{(\inverse\csc(x))}{x}$ is always negative.
%   \item What is the domain of $\inverse\sec(x)$?
%   \end{enumerate}
% \end{embeddedproblem-enumerate}
We now have  differentiation formulas for all of the inverse
trigonometric functions. These, along with the constraints on their
ranges are given in the table below.

\alert{Memorize them!} You will need them.

      %       \begin{table}[h]
      %       \caption{Differentials of Inverse Trigonometric functions}
      %       \label{tab:InvTrigDiff2}
\begin{center}
  \begin{tabular}{|cccc|} \hline Function\amp{}
                                           Domain\amp{}Range\amp{}Derivative\\\hline\hline $y=\inverse\tan(x)$ \amp{} Any
                                                                                                       real number \amp{} $-\frac{\pi}{2}\lt{}y\lt{}\frac{\pi}{2}$
                                                       \amp{}$\displaystyle\dfdx{y}{x}=\frac{1}{1+x^2}$\\
                                         \amp{}\amp{}\amp{}\\
    $y=\inverse\cot(x)$ \amp{} Any real number \amp{} $0\lt{}y\lt{}\pi$
                                                       \amp{}$\displaystyle\dfdx{y}{x}=\frac{-1}{1+x^2}$\\
                                         \amp{}\amp{}\amp{}\\
    $y=\inverse\sin(x)$ \amp{} $-1\le x\le 1$ \amp{}
                                           $-\frac{\pi}{2}\le
                                           y\le\frac{\pi}{2}$
                                                       \amp{}$\displaystyle\dfdx{y}{x}=\frac{1}{\sqrt{1-x^2}}$\\
                                         \amp{}\amp{}\amp{}\\
    $y=\inverse\cos(x)$ \amp{} $-1\le x\le 1$ \amp{} $0\le y\le\pi$
                                                       \amp{}$\displaystyle\dfdx{y}{x}=\frac{-1}{\sqrt{1-x^2}}$\\
                                         \amp{}\amp{}\amp{}\\
    \multirow{2}{*} {$y=\inverse\sec(x)$} \amp{}$x\le-1$
                                                 \amp{}$0\le y\lt{}\frac{\pi}{2}$ \amp{}\multirow{2}{*}
                                                                           {$\displaystyle\dfdx{y}{x}=\frac{1}{\abs{x}\sqrt{x^2-1}}$}\\
                                         \amp{} or $ x\ge 1$
                                                 \amp{}or $\frac{\pi}{2}\lt{}y\le\pi$\amp{}\\
                                         \amp{}\amp{}\amp{}\\
    \multirow{2}{*}{$y=\inverse\csc(x)$} \amp{}$x\le-1$
                                                 \amp{}$0\le y\lt{}\frac{\pi}{2}$
                                                       \amp{}\multirow{2}{*}{$\displaystyle\dfdx{y}{x}=\frac{-1}{\abs{x}\sqrt{x^2-1}}$}\\
                                         \amp{} or $ x\ge 1$
                                                 \amp{}or $-\frac{\pi}{2}\lt{}y\le0$\amp{}\\
                                         \amp{}\amp{}\amp{}\\\hline
  \end{tabular}
\end{center}
      %       \end{table}

Until now there has been a pattern to differentiation that we have not
remarked on. But you may have noticed that when we differentiate a
function we usually get back another function of the same type.
That is,
\begin{enumerate}
\item The derivative of a polynomial is another polynomial,
\item The derivative of a quotient of polynomials is another quotient
  of  polynomials,
\item The derivative of a trig function is another trig
  function.
\end{enumerate}

The inverse trigonometric functions break this pattern. The derivative
of an inverse trigonometric function is not another inverse
trigonometric function. Nor is it a trig function, a polynomial, or a
quotient of polynomials\aside{Well, actually that's not quite true the
  derivative of the inverse tangent \emph{is} a quotient of
  polynomials, and a particularly simple one at that.}.

Do you suppose there is any significance to this?





      %       when package did not graph
      %       $y=\arctan\left(\frac{1+x}{1-x}\right)$, but actually graphed
      %       $y=\inverse\tan\left(\frac{1+x}{1-x}\right)$.  So, what is the
      %       difference?  This brings up what we actually mean by a function.  In
      %       truth, $y=y(x)=\arctan(x)$ is not really a function.  To be a
      %       function, then for every $y$ value, there must be a single value that
      %       we can call $y(x)$. 

      %       Notice that this is \emph{not} the graph of a function.  This is
      %       because for any integer, $n$, $x=\tan(y)=\tan(x+n\pi)$.  
      %       Because of this restriction on the \emph{range} of $\tan(x)$ (and thus
      %       one the \emph{domain} of $\inverse\tan(x)$), when    $x\lt{}1$
      %       $\arctan\left(\frac{1+x}{1-x}\right)=\arctan(x)+\frac{\pi}{4}$.
      %       However, for $x\gt{}1$, $\arctan\left(\frac{1+x}{1-x}\right)=\arctan(x)-\frac{3\pi}{4}$.

      %       This can be clearly seen when we superimpose the graphs of
      %       $\arctan(x)+\frac{\pi}{4}$ (in orange) and
      %       $\arctan(x)-\frac{3\pi}{4}$ (in red)  on the two branches of
      %       $\arctan\left(\frac{1+x}{1-x}\right)$ (in black) as in the following sketch. \\
      %       \centerline{\includegraphics*[height=1.5in,width=4in]{../Figures/arctan4}}


      %       An easy way to remember what $\inverse{\tan}(x)$ means is to read
      %       the symbol $\inverse{\tan}(x)$ as ``the angle whose tangent is
      %       $x.$'' This emphasizes that $x$ is the tangent of some angle, and
      %       $\inverse{\tan}(x)$ is that angle.

      %       There similar issues with the ranges of the other five inverse
      %       trigonometric functions, $\inverse\sin,$ $\inverse\cos,$
      %       $\inverse\cot,$ $\inverse\sec,$ and $\inverse\csc$, but it is harder
      %       to see how to resolve them than it was for the inverse
      %       tangent. Because these problems consist mostly of defining
      %       consistent conventions for these functions we will just tell you
      %       what conventions are rather than discussing the conventions
      %       themselves at length as we have here.

      %       \begin{ProficiencyDrill}
      %       Because $\sec(x)=\frac{1}{\cos(x)}$ the
      %       restriction on the \emph{range} of the inverse secant function (and on
      %       the \emph{domain} of the secant) is the same as the restriction on the
      %       cosine and its inverse.

      %       Use this to show that
      %       $$
      %       \dx\left(\arcsec(x)\right)=\frac{1}{\abs{x}\sqrt{x^2-1}}.
      %       $$
      %       \hint{Don't let this problem frighten you. The absolute value
      %       follows directly from the \emph{domain} restriction. Use a triangle in
      %       the unit circle as before to get the formula.}
      %       \end{ProficiencyDrill}

%\pagebreak[4]
\begin{mynotation}{Inverse Function}

      %       It is easy to understand what the inverse of a function is: If the
      %       function $\tan$ takes $x$ to $y$, then the function $\inverse\tan$
      %       ``undoes'' the function $\tan$ (If $y=\tan(x)$ then
      %       $x=\inverse\tan(y)$).

      %       The \emph{concept} really is that simple.


      %       But\notefrombob{I have also seen the notation y=Arctan x to denote the
      %       principal branch instead of y=x .  Do we want to utilize this at all
      %       in the discussion about inverse functions or is this complicating
      %       things unnecessarily? -- Jul 19, 2021} function inversion can be
      %       \emph{very} difficult to work with. This is partly because it is so
      %       abstract, but also because the notation we use can be very
      %       confusing. Probably the clearest notation we have is the ``arc''
      %       designation. The symbol ``$\arctan,$'' for example, clearly signals
      %       that the function\aside{In the computer age this has been shortened to
      %       ${\rm atan}$ for reasons that are no longer apparent.}, $\arctan(x)$
      %       returns an angle (the ``arc'' of a circle).


      %       We (the authors) have consistently used the ``arc'' prefix to
      %       designate the multifunctions which gave us the inverse trigonometric
      %       functions. Thus by our convention $\arctan(x)$ and $\inverse\tan(x)$ refer
      %       to different concepts altogether. We will continue this practice throughout this
      %       text. That is, we  will always, and only, refer to inverse functions by placing
      %       ``$-1$'' as an exponent on the name of the original function. Thus if
      %       $f(x)$ is an invertible function then $\inverse f(x)$ is it's inverse.

      %       We do this partly to avoid confusion in what can be a very confusing
      %       topic. But this is by no means a standard usage. 
      %       In fact for most scientists, engineers, and mathematicians
      %       ``$\arctan(x)$'' and $\inverse\tan(x)$, are synonymous.  In the
      %       computer age it has also become common to use $\text{atan}(x)$, and
      %       there are other variants as well.


  The use of ``$-1$'' as an exponent is probably the most common
  notation used to indicate function inversion. Unfortunately, from the
  standpoint of a student it is also probably the worst possible
  notation we could have invented because it is so very easily confused
  with reciprocation. While it \underline{\alert{is true}} that
  $$2^{-1} = \frac12,$$  it is \underline{\alert{not true}} that
  $$\inverse{\sin}(x) = \frac{1}{\sin(x)}.$$

  \begin{wrapfigure}[19]{r}{3.2in}
%    \vskip-6mm{}
    \captionsetup{labelformat=empty}
%    \centerline{\href{https://mathwithbaddrawings.com/}{\includegraphics*[height=1.5in,width=2.2in]{../Figures/BadDrawing-InvSine}}}
    \centerline{\href{https://mathwithbaddrawings.com/}{\includegraphics*[height=2.5in,width=3.2in]{../Figures/BadDrawing-InvSine}}}
    \label{fig:BadDrawing-InvSine}
    \caption{\textcolor{black}{\centerline{The World's Most Evil Mathematician}
        \centerline{from {Math
            With Bad Drawings}}
          \centerline{by Ben Orlin}}}
  \end{wrapfigure}
  That\aside{We've made the ``World's Most Evil Mathematician'' even
    more evil by defining $\arcsin(x)$ to be a multifunction. This is
    actually a conceit of the authors of this text. Most of the world
    uses $\arctan(x)$ and $\inverse\tan(x)$
    interchangeably. Eventually you will too, but for now we're hoping
    this will aid your efforts to learn.} the ``$-1$'' notation is
  used for both comes from the fact that both $2^{-1}$ and
  $\inverse\sin(x)$ really are inverses. But they are different kinds
  of inverses: $\inverse2$ is the \emph{multiplicative} inverse of
  $2$, which means simply that
  $$\frac12\cdot2=1 \text{ and that } 2\cdot\frac12 =1.$$ 
  On the other hand $\inverse\sin(x)$ is the \emph{functional} inverse of
  $x=\sin(y)$, which means that
  $$
  \inverse\sin(\sin(x)) = x, \text{(if $-\frac{\pi}{2}\lt{}x\lt{}\frac{\pi}{2}$)}
  $$
  and that
  $$
  \sin(\inverse\sin(x))=x.
  $$

  In the former case you can think of the ``$-1$'' in the exponent as
  an operator. It ``operates on'' $2$ by taking its reciprocal. In the
  latter case the ``$-1$'' is part of the name of the function. If you
  read $\inverse\sin(x)$  as ``the angle, between $-\pi/2$ and $\pi/2$ whose sine is''
  it is a little less confusing. Similarly for the other inverse
  trigonometric functions.
\end{mynotation}

The inverse trigonometric functions will not have a significant role
in our story for a while. However they will come to be very useful in
some topics that will come up next semester. We've introduced them
here simply because it makes sense to develop all of the
differentiation rules at more or less the same time.

%\pagebreak[4]
\begin{ProficiencyDrill}
  Compute each of the following:
  \begin{enumerate}[label={\bf{}(\alph*)}]
    \begin{multicols}{2}
    \item $\dx(\inverse\sin(x^{1/2})$
    \item $\dx(\inverse\tan(x^{1/2})$
    \item $\dx(\inverse\sec(x^{1/2})$
    \item $\dx \left( \inverse\sin(x) \right)^{-1}$
    \item $\dx \left( \inverse\tan(x) \right)^{-1}$
    \item $\dx \left( \inverse\sec(x) \right)^{-1}$
    % \item
    %   $\dx \left[ (3x^2+2x-1) \inverse\sec(\inverse{x}+x^2) \right]$
    % \item $\dx \left(\inverse\csc(x^2) \right)$
    \end{multicols}
  \end{enumerate}
\end{ProficiencyDrill}

\pagebreak[4]
\begin{embeddedproblem}[Find the Pattern]
  Show that for any constant $a\gt{}0$,
  \begin{enumerate}[label={\bf{}(\alph*)}]
    \begin{multicols}{2}
    \item
      $\dx{\left(\dfrac1a\inverse\tan\left(\dfrac{x}{a}\right)\right)}=\dfrac{1}{x^2+a^2}\dx{x}$
    \item
      $\dx{\left(\dfrac1a\inverse\sin\left(\dfrac{x}{a}\right)\right)}=\dfrac{1}{a\sqrt{a^2-x^2}}\dx{x}$
    \item
      $\dx{\left(\dfrac1a\inverse\sec\left(\dfrac{x}{a}\right)\right)}=\dfrac{1}{\abs{x}\sqrt{x^2-a^2}}\dx{x}$
    \item
      $\dx{\left(\dfrac1a\inverse\cot\left(\dfrac{a}{x}\right)\right)}=\dfrac{1}{x^2+a^2}\dx{x}$
    \item
      $\dx{\left(\dfrac1a\inverse\cos\left(\dfrac{a}{x}\right)\right)}=\dfrac{1}{\abs{x}\sqrt{x^2-a^2}}\dx{x}$
    \item
      $\dx{\left(\dfrac1a\inverse\csc\left(\dfrac{a}{x}\right)\right)}=\dfrac{1}{a\sqrt{a^2-x^2}}\dx{x}$
   \end{multicols}
 \end{enumerate}
 Identify and describe any patterns you see in the forms of these derivatives.
\end{embeddedproblem}

\begin{embeddedproblem}
  Show that
  $\displaystyle\dx{\left(\dfrac{\inverse\sin(x)+x\sqrt{1-x^2}}{2}\right)}=\sqrt{1-x^2}\dx{x}$.
\end{embeddedproblem}

\begin{embeddedproblem-enumerate}[Find the Pattern]
  \begin{enumerate}[label={\bf{}(\alph*)}]
    \item   Show that each of the following statements is true.
  \begin{enumerate}[label={\bf{}(\roman*)}]
    \begin{multicols}{2}
    \item
      $\dx\left(\inverse\sin\left(\frac{x}{\sqrt{x^2+1}}\right)\right)=\frac{1}{\sqrt{x^2+1}}\dx{x}$
    \item
      $\dx\left(\inverse\sin\left(\frac{x}{\sqrt{x^2+4}}\right)\right)=\frac{2}{\sqrt{x^2+4}}\dx{x}$
    \item
      $\dx\left(\inverse\sin\left(\frac{x}{\sqrt{x^2+9}}\right)\right)=\frac{3}{\sqrt{x^2+9}}\dx{x}$
    \item
      $\dx\left(\inverse\sin\left(\frac{x}{\sqrt{x^2+5}}\right)\right)=\frac{\sqrt{5}}{\sqrt{x^2+5}}\dx{x}$
    \end{multicols}
  \end{enumerate}
  \item Now compute
    $\dx\left(\inverse\sin\left(\frac{x}{\sqrt{x^2+a^2}}\right)\right),$
    where $a\in\RR.$
  \item  Does this problem change if $a$ is negative? How?
  \end{enumerate}
  \vskip-4mm{}
\end{embeddedproblem-enumerate}


\pagebreak[4]
\begin{embeddedproblem}
  Show that each of the following statements is true.
  \begin{enumerate}[label={\bf{}(\alph*)}]
%    \begin{multicols}{2}
    \item
      $\dx\left(x\inverse\sin(x)+\sqrt{1-x^2}\right) =
      \inverse\sin(x)\dx{x}.$
    \item
      $\dx\left(x\inverse\cos(x)-\sqrt{1-x^2}\right) =
      \inverse\cos(x)\dx{x}.$
 %   \end{multicols}
    \end{enumerate}
    \vskip-6mm{}
\end{embeddedproblem}

%\pagebreak[4]
\begin{embeddedproblem-enumerate}
  % Consider $y_1(x)=\inverse\cot\left(\frac{1-x}{1+x}\right),$
  % and $y_2(x)=\inverse\tan(x).$
  \begin{enumerate}[label={\bf{}(\alph*)}]
  \item Show that
    $\dx\left(\inverse\cot\left(\frac{1-x}{1+x}\right)\right)=
    \dx\left(\inverse\tan(x)\right)$.
  \item Part (a) says that the graphs of
    $y=\inverse\cot\left(\frac{1-x}{1+x}\right)$ and
    $y=\inverse\tan(x)$ have the same slope when $x\neq-1$  and that
    their graphs  should differ by a constant. Substitute $x=0$
    into both to see what this constant should be.
  \item Now plot the graphs of both. Do they differ by a constant?
  \item Can you explain this? \comment{You might want to review
      Digression~\#\ref{digression:dzEqZero} before you answer.}
    % \begin{multicols}{2}
  % \item Graph $y_1(x)$ and $y_2(x)$ on the same set of axes. You
  %   should see that when $x$ is positive they have \emph{nearly} the
  %   same graph.
  % \item Show that $\dx{y_1}$ and $\dx{y_2}$ are also \emph{nearly}
  %   the same, at least when $x\ge0.$
  % \item Modify $y_2$ so that it is identical to $y_1.$
    % \end{multicols}
  \end{enumerate}
  \vskip-6mm{}
\end{embeddedproblem-enumerate}

\section{Curvature}
\label{sec:curvature}
\TLogo{appendix:radi-curv-accord}

\begin{wrapfigure}[]{r}{2.7in}
  \vskip-6mm{}
  \captionsetup{labelformat=empty}
  \centerline{\includegraphics*[height=2.5in,width=2.7in]{../Figures/Curvature6}}
  \label{fig:Curvature6}
\end{wrapfigure}
Intuitively speaking, the \term{curvature} of a plane curve is how
quickly a curve turns. It is an important property to understand when
designing, for example, a road or a roller-coaster where the sharpness
of a curve determines the safe speed limit for that curve. Also, the
stresses imposed on the physical structure of a moving object during
any high speed maneuver are directly related to the curvature of their
paths.

One of the first people to study the notion of curvature was \href{https://mathshistory.st-andrews.ac.uk/Biographies/Oresme/}{Nicole
  Oresme} ($1323-1382$). He concluded the curvature of a straight line
should be zero (no surprise) and that the curvature of a circle should
be inversely proportional to the radius of the circle. For example, a
circle of radius $1$ has twice as much curvature as a circle of radius
$2$ and the radius of the earth is so large that the curvature of say,
the equator, is imperceptible near the surface.

Because he was the first, Oresme's approach to curvature
qualitative. We'd like to be more precise. Specifically, we'd like to
have a number we can associate with each point on the curve which
tells us the curvature at that point. Ideally, this would recapture
Oresme's conclusions that a straight line should have a curvature of
zero at every point and that the curvature of a circle should be be
constantly equal to the reciprocal of the radius. The arctangent
(multi)function is exactly the tool we need to begin.

Consider the circles in the sketch above.
Notice that on the unit circle, it takes $\frac{\pi}{2}$ units to turn
the red tangent lines
$\frac{\pi}{2}$ radians. On the circle of radius $2$, it takes
$2 \times \frac{\pi}{2}$ units to turn the tangent the same $\frac{\pi}{2}$
radians, whereas on the circle of radius $1/3$, it takes
$\frac13\times\frac{\pi}{2}$ units to turn it the same $\frac\pi2$
radians.  This suggests that a possible definition for
\term{curvature} might be the number of radians the tangent line
rotates as it turns divided by the distance traveled while making the
turn. In other words, the curvature for a circle of radius $r$ would
be
$$
\frac{\theta \text{ radians }}{r\theta \text{ units}}=\frac{1}{r}
\text{ radians per unit. }
$$
This seems to make intuitive sense too. A circle with a radius of $0.0000001$
meters certainly seems to be ``more curved'' than a circle with radius
$1,000,000,000$ meters.

By the same reasoning, the curvature of a straight line would be
zero. (Why?) 

This definition works fine for circles and straight lines which have
constant curvature, but what about a more general curve. What if the
curvature changes from point to point along the curve. Think of a winding
road.

Clearly we'll need a definition of curvature that also changes from point
to point on the curve so it seems reasonable to consider infinitely
small changes in both the angle $\theta$ and in arc length $s$.  More
precisely, suppose we have a point $P$ with coordinates $(x,y)$ lying
on a curve $C$.


\begin{wrapfigure}[]{l}{1.5in}
  \vskip-6mm{}
  \captionsetup{labelformat=empty}
\centerline{\includegraphics*[height=1.5in,width=1.5in]{../Figures/Curvature5}}
\label{fig:Curvature5}
\end{wrapfigure}
If we let $\theta$ represent the angle formed in the figure at the
left by the tangent line to the curve $C$, at point $P$, with some
fixed line
then the curvature is defined to be
\begin{equation}
\kappa=\abs{\dfdx{\theta}{s}}.\label{eq:CurvatureDef}
\end{equation}
To keep things simple we'll always measure the angle from a horizontal
line. 

It can be difficult at first to see how this formula measures
curvature but think carefully for a moment about the meaning of the
expression $\dfdx{\theta}{s}$. Loosely put ``the rate of change of
$\theta$ (the angle of our curve with the horizontal) with respect to
$s$ (the arclength of our curve)'' is the ratio of how much a point
moving on  the curve
has changed direction to how far it has moved. The more our curve
changes direction in a given distance the more it is curved; the more
curvature it has.


Don't let the absolute value bars in this formula disturb you. They
are there simply because we don’t want to consider whether a curve
bends to the right or to the left (whether $\theta$ is increasing or
decreasing), but just how fast it bends.

Also note, for later reference, that
\begin{equation}
  \tan(\theta)=\dfdx{y}{x}.\label{eq:TanEqDeriv}
\end{equation}

% \begin{ProficiencyDrill}
%   Do you see that this definition captures the idea that curvature is
%   a measure of how quickly, or slowly, our direction changes as we
%   travel along the curve?  Explain.
% \end{ProficiencyDrill}

A good test of the value of a new definition is whether or not it
produces the same results that our (actually Oresme's) earlier, more
intuitive musings did. It should be clear that with this definition,
the curvature of a straight line would still be zero. (Why?)

\begin{embeddedproblem}
  \label{prob:CircleCurvature}
  Use the following picture of a quarter-circle with radius $r$ to show that our
  general definition of curvature yields $\kappa =\frac{1}{r}$.\\
  \centerline{\includegraphics*[height=1.6in,width=1.9in]{../Figures/Curvature3}}
\end{embeddedproblem}

In problem (\ref{prob:CircleCurvature}) it was helpful to know that the
line tangent to a circle is always perpendicular to its radius. But
for an arbitrary curve this is not necessarily true. Let's look at the
general situation.


First recall from equation~(\ref{eq:TanEqDeriv}) that
$ \tan(\theta) = \dfdx{y}{x}$.  Rewriting this as
$\theta=\arctan\left(\dfdx{y}{x}\right)$ and differentiating we have
\begin{equation}
  \dx{\theta}=\frac{1}{1+\left(\dfdx{y}{x}\right)^2}\dx\left(\dfdx{y}{x}\right).\label{eq:DiffTheta}
\end{equation}

To continue we need to compute $\dx\left(\dfdx{y}{x}\right)$ which
seems straightforward enough. Since $\dfdx{y}{x}$ is a quotient the
Quotient Rule applies. This results in
$$
\dx\left(\dfdx{y}{x}\right) =
\frac{\dx{x}\textcolor{red}{\dx{(\dx{y})}}-\dx{y}\textcolor{red}{\dx{(\dx{x})}}}{(\dx{x})^2}.
$$
Next we simplify the notation a bit. If we let\aside{Note this very
  carefully, $\d{(\dx{y})}\neq (\dx{y})^2$,  because  the latter is
  $(\dx{y})\times(\dx{y})$ whereas the former is the \emph{second
    difference} of $y$.} $\textcolor{red}{\d{(\dx{y})}}=\dx^2{y}$ and
$\textcolor{red}{\d{(\dx{x})}}=\dx^2{x}$ then we have
$$
\dx\left(\dfdx{y}{x}\right) =
\frac{\dx{x}\d^2{y}-\dx{y}\dx^2{x}}{(\dx{x})^2}.
$$
which seems OK at first glance. But there is a problem because
$\dx^2x$ is equal to zero, but $\dx^2y$ is \emph{not necessarily}
equal to zero and it is not at all clear why this should be true. We
will take a moment to investigate this.

\pagebreak[4]
\begin{Digression}[\foreign{Hic Sunt Dracones} (Here Be Dragons)]
  \label{digression:curvature-1}
\begin{wrapfigure}[13]{r}{1.5in}
  \vskip-4mm{}
\captionsetup{labelformat=empty}
\centerline{\href{https://en.wikipedia.org/wiki/Hunt-Lenox_Globe}{\includegraphics*[height=1.5in,width=1.5in]{../Figures/LennoxGlobe2}}}
\caption{From the  Rare Book Division of the The New York Public Library}
\label{fig:LennoxGlobe}
\end{wrapfigure}
In medieval times mapmakers would fill in the unexplored regions of
their maps with illustrations of dragons and other mythological beasts
to indicate that these uncharted waters were the most dangerous. The
Lennox Globe (1510 CE, seen at the right) even has the inscription
``\foreign{Hic Sunt Dracones}'' (Here Be Dragons) along the eastern
Asian coast. (See the inset in the figure at the right.)
  
The concept of the differential of a differential $(\dx^2x)$ is a
mathematical dragon of sorts. If we follow it too far we risk bringing
all of  Calculus crashing to the ground in a fiery storm of logical
contradictions.

  We will proceed carefully.

   The problem here is conceptual, not computational. That is, there is
  no computation we can do to see why   $\dx^2x$ is equal to
  zero. Rather we need to think carefully about the meanings of our
  symbols.

  \subsubsection*{Leibniz' Conception}
  When we use Leibniz' differential notation we think of $\dx{y}$
  and $\dx{x}$ as tiny increments in the $x$ and $y$
  coordinates. Specifically, $\dx{x}=x_2-x_1$ where $x_2$ and $x_1$
  are \emph{infinitely} close together. As we've seen the concept of
  a differential is already a difficult thing to accept.

  But if we stick with Leibniz then it must be that
  $ \dx(\dx{x}) = \dx{x_2}-\dx{x_1}, $ right? We can write down, and
  read, 
  the symbols but what they seem to \emph{mean} is that $\dx(\dx{x})$
  is the infinitely small difference of two infinitely small
  differences.

  \vskip2mm{} \centerline{\bf{} Down that path lies madness. We won't
    go there.}  \vskip2mm{}


  \subsubsection*{Newton's Conception}
  Instead we'll adopt Newton's viewpoint and hope that it gives us a
  meaningful interpretation of our symbols.
  
  For Newton time was the only variable, and time always flows forward
  at a constant rate\aside{Of course, with Einstein's work on
    Relativity Theory we now know that this assumption is not true. No
    matter. Newton assumed it was true and he based his Method of
    Fluxions on that assumption.}. He thought of a variable, $x$ for
  example, as moving, or flowing (\foreign{fluent}) in time and
  $\dot{x}$ was the rate of flow (\foreign{fluxion}) of the
  \foreign{fluent} $x$. In the simplest example if $x$ is the position
  of a point then $\dot{x}$ is its velocity.

  Newton only applied his dot notation to something that was changing
  in time; something that has a \emph{rate of change with respect to
    time.} Can we apply it to time itself?

  Sure we can. Time itself is ``changing in time'', isn't it? Time
  flows at a rate of $1$ second per second, or one day per day, or one
  century per century. The units don't really matter. Loosely speaking
  (and mixing the Newtonian and Leibnizian viewpoints a bit), an
  increment of time now has the same magnitude as an increment of time
  later, has the same magnitude as an increment of time in the
  past. That is, every increment of $t$ is the same size.  If we let
  $t$ represent elapsed time then this means that the infinitesimal
  increment $\dx{t}$ is \emph{constant}, so the fluxion of $\dx{t}$ is
  zero: $\dx(\dx{t})=\dx^2{t}=\dot{\dx{t}}=0$.  This is not true of a quantity that is
  flowing at a variable rate. For example, if $y=t^2$ then
  $\eval{\dx{y}}{t}{1}=2\dx{t}$ whereas $\eval{\dx{y}}{t}{2}=4\dx{t}$.

  This solves our problem. If $\dx{t}$ is \emph{constant} then by the
  Constant Rule $\dx^2{t}=\dx(\dx{t})=0$.

  Moreover we always have the option of adopting Newton's convention
  that the \emph{only} independent variable is time so that when we
  write $\dfdx{y}{x}$ we are thinking of  the independent
  variable $x$ as time (although we rarely say so out loud), and we are
  thinking of $y$ as functionally dependent on time (represented by
  $x$ in this case).  Therefore
  $$\dx^2x=\dx(\dx{x})=0 \text{ and } \dx^2y=\dx(\dx{y})\neq0.$$

  % Recall that the differential notation was Leibniz' invention and it
  % reflects a static viewpoint. But Newton's approach was dynamic. For
  % Newton the only \alert{independent} variable was time. Everything
  % else depends on how much time has elapse since we started our
  % clock. That is, everything else is a dependent variable (a
  % \term{function} of time).

  % Moreover Newton assumed that time flows forward into the future at a
  % \alert{constant} rate\aside{Of course, with Einstein's work on
  % Relativity Theory we now know that this assumption is not true.}
  % Since the differential of a
  % constant is zero we have, by the Constant Rule $\dx{(\dx{t})}=0$.


  As a practical matter we will rarely be interested in second order
  differentials like $\dx^2x$. As we indicated in
  section~\ref{sec:more-high-deriv} differentials in general are
  really just an aid to intuition and a step along the way to finding
  derivatives (differential ratios). Viewed in this light Leibniz'
  notation is again very helpful for if $y=f(x)$ then the second
  derivative of $y$ is
  $$
  \dfdx{}{x}\left(\dfdx{y}{x}\right) = \dfdxn{y}{x}{2}
  $$
  as before.

  But now there is another difficulty that can't be ignored.

  In order to make sense of the expression $\dx^2{x}$ we had to
  abandon Leibniz' conception altogether. It is nice that Newton's
  viewpoint gives us a way to make sense of the notation but what this
  tells us is that taking either viewpoint (Newton's or Leibniz') does
  not get us to the logical heart of the matter. A viewpoint is just a
  viewpoint. It is a manner of thinking. We adopt a viewpoint in
  order to build intuition.

  When two viewpoints are at odds like this it means that we need
  something more general. In this case it means we need a theory that
  subsumes, and thus replaces, the views of both Leibniz and Newton.

  Eventually we will need to stand at the top of the Calculus mountain
  so we can see it all at once. From that vantage point we'll be able
  to see the path that Newton took, and the path that Leibniz took.
  This issue must be addressed but for now we will continue to work
  intuitively.  We will begin our own trek to the summit in
  chapter~\ref{chapt:whats-wrong-with}. .
\end{Digression}


Returning to our computation in progress we have
\begin{align*}
  \dx\left(\dfdx{y}{x}\right) \amp{}=
                                \frac{\dx{x}\cdot
                                \d^2{y}-\dx{y}\cdot\CancelToRed{0}{\d^2{x}}}{(\dx{x})^2}\\
  \dx\left(\dfdx{y}{x}\right) \amp{}= \left(\dfdxn{y}{x}{2}\right)\dx{x}% \frac{\dx{x}\cdot\d^2{y}}{(\dx{x})^2}
\end{align*}
Returning to equation~(\ref{eq:DiffTheta}) we have
\begin{align*}
  \dx{\theta}\amp{}=\frac{1}{1+\left(\dfdx{y}{x}\right)^2}\dx\left(\dfdx{y}{x}\right)\\
             \amp{}=\frac{1}{1+\left(\dfdx{y}{x}\right)^2}\left(\dfdxn{y}{x}{2}\right)\dx{x}% \frac{\dx{x}\cdot\d^2{y}}{(\dx{x})^2}\\
\end{align*}


\begin{embeddedproblem}
  Use the fact that $\dx{s}=\sqrt{(\dx{x})^2+(\dx{y})^2}$, to show that
  \begin{equation}
    \kappa=\frac{\abs{\dfdxn{y}{x}{2}}}{\left[1+\left(\dfdx{y}{x}\right)^2\right]^{\frac32}}.\label{eq:CurvatureDef2}
  \end{equation}
  \vskip-6mm{}
\end{embeddedproblem}

Since the equation of a straight line is $y=mx+b$ it is clear that
formula~(\ref{eq:CurvatureDef2}) recovers Oresme's assertion that the
curvature of a straight line must be zero. (Why?)

Let's apply formula~(\ref{eq:CurvatureDef2}) to a circle to see if this
agrees with Oresme's definition of the curvature of a circle. 
Consider the circle of radius $r$ given by the graph of $x^2+y^2=r^2$.
Differentiating, we obtain $  2x\dx{x}+2y\dx{y}=0$ so that \\
\centerline{$\displaystyle \dfdx{y}{x}=-\frac{x}{y}.  $}
Differentiating both sides of this gives\\
%$\displaystyle\dx\left(\dfdx{y}{x}\right)=\dfdxn{y}{x}{2}=\frac{-y\dx{x}+x\dx{y}}{y^2}
%=\frac{-y+x\dfdx{y}{x}}{y^2}$.
\centerline{$\displaystyle
  \dx\left(\dfdx{y}{x}\right)=\frac{-y\dx{x}+x\dx{y}}{y^2}$}
or, after dividing both sides by $\dx{x}$\\
\centerline{$\displaystyle  \dfdxn{y}{x}{2}=\frac{-y+x\dfdx{y}{x}}{y^2}$.}



% \begin{ProficiencyDrill-enumerate}
%     \begin{enumerate}[label={\bf{}(\alph*)}]
%     \item   Use the  values for $\dfdx{y}{x}$ and $\dfdxn{y}{x}{2}$
%       just computed to show that for a circle of radius $r$, the
%       curvature $\kappa=1/r$.
%     \item Does this agree with the idea that as the radius of a circle
%       gets larger, the curvature gets smaller?
% \end{enumerate}
% \end{ProficiencyDrill-enumerate}

\begin{ProficiencyDrill}
   Use the  values for $\dfdx{y}{x}$ and $\dfdxn{y}{x}{2}$ we
      just computed to show that for a circle of radius $r$, the
      curvature $\kappa=1/r$.
 Does this agree with the idea that as the radius of a circle
      gets larger, the curvature gets smaller?
    \end{ProficiencyDrill}
    

As we said earlier, the curvature of a general curve need not be
constant, but we can still apply formula~(\ref{eq:CurvatureDef2}).

\begin{embeddedproblem}
  Show that the curvature of the parabola $y=x^2$ is given by
  $
  \displaystyle\kappa=\frac{2}{(1+4x^2)^{\frac32}}$.
Where does the greatest curvature occur? Does this agree with what
  you can see on the graph of $y=x^2$?
\end{embeddedproblem}

\begin{embeddedproblem}
Consider the ellipse given by $\frac{x^2}{a^2}+\frac{y^2}{b^2}=1$,
  $a\gt{}b\gt{}0$.\\
%   \begin{wrapfigure}[5]{r}{2.3in}
%   \vskip-21mm{}
% \captionsetup{labelformat=empty}
  \centerline{\centerline{\includegraphics*[height=1.4in,width=2.3in]{../Figures/Curvature7}}}
% \label{fig:Curvature7}
% \end{wrapfigure}
  \begin{enumerate}[label={\bf{}(\alph*)}]
  \item Show that the curvature is given by
    $ \kappa=\frac{a^4b^4}{(a^4y^2+b^4x^2)^{\frac32}}$.\\
    \hint{You may want to rewrite
      $\frac{x^2}{a^2}+\frac{y^2}{b^2}=1$ as
      $b^2x^2+a^2y^2=a^2b^2$.
    }
  \item According to this graph, the curvature at $(a,0)$ should be
    greater than at $(0,b)$. Use the formula you just derived to
    verify this.
  \item A circle is the special case of an ellipse when $a=b$. Use the
    formula from part (a) to compute the curvature of a circle.
  \end{enumerate}
  \vskip-6mm{}
\end{embeddedproblem}
%\vskip-3mm{}

\begin{embeddedproblem}[Descartes' Method of Normals and Curvature]
  Recall that in exercise~\#\ref{prob:DescartesNormals2} of
  section~\ref{sec:descartes-normals} we used Descartes's Method of
  Normals to find the
    \begin{wrapfigure}[7]{r}{2.4in}
%    \vskip-4mm{}
    \captionsetup{labelformat=empty}
    % \centerline{\includegraphics*[height=1.5in,width=2in]{../Figures/DescartesCircle5}}
    \centerline{\includegraphics*[height=1.8in,width=2.4in]{../Figures/DescartesCircle5}}
    \label{fig:DescartesCircle3}
  \end{wrapfigure}
line tangent to the graph of $y=\sqrt{2x}$ at
  the point $(2,2)$ by finding a circle with its center on the
  $x$-axis which touches the graph exactly once, as seen in the 
  sketch at the right.
  \begin{enumerate}[label={\bf{}(\alph*)}]
  % \item Take a guess, do you think the circle and the parabola in this
  %   sketch have the same curvature at $(2,2)$?
  \item At their point of intersection compute the curvatures of the
    parabola and the circle seen in the diagram at the right and show
    that they are not the same.
  \item Find the equation of the circle which passes through the point
    $(2,2)$, and also has  the same slope
    and curvature as $y=\sqrt{2x}$ at the point $(2,2)$. This is known as the \term{osculating circle},
      or ``kissing circle''.
  \item Show that the center of the osculating circle is on the line
    which passes through $(2,2)$ and $(3,0)$.
  \end{enumerate}
  \vskip-6mm{}
\end{embeddedproblem}

We observed earlier that if we  define curvature using
equation~(\ref{eq:CurvatureDef}) then the curvature of a straight line
will be zero as Oresme said it should be. The converse is also
true. Specifically, if $\abs{\dfdx{\theta}{s}}=\kappa=0$ then $\theta$
must be a constant and so $\dfdx{y}{x}=\tan(\theta)$ must also be a
constant. (For simplicity we'll ignore the case where
$\theta=\pm\frac{\pi}{2}$.)

\begin{ProficiencyDrill}
  Assume that  $-\frac{\pi}{2}\lt \theta \lt
  \frac{\pi}{2}$ and that $\theta$ is constant. Complete the line of
  reasoning we began in the previous paragraph to show that
  $$
  y=  \left(\tan(\theta)\right)x+b
  $$
  for some number $b$. Explain how you know that this is this the
  equation of a line.
\end{ProficiencyDrill}

\begin{wrapfigure}[]{r}{2in}
\captionsetup{labelformat=empty}
\centerline{\includegraphics*[height=1.5in,width=1.5in]{../Figures/Curvature5}}
\label{fig:Curvature5Redux}
\end{wrapfigure}We've also shown that the curvature of a circle is
also constant (specifically, it is the reciprocal of the circle's
radius). The converse of this is also true\aside{But only if we
  confine our attention to curves in the plane. It turns out that
  there are many curves in three-dimensional space, (that is, curves
  that are not constrained to a two-dimensional plane) with constant
  curvature which are neither circles nor lines.}. This is shown in
much the same manner as the straight line but, naturally, the
computations are more complex.

To see this recall the diagram that we used to define curvature, shown
at the right. If we  suppose that $\abs{\dfdx{\theta}{s}}=\kappa$ then
we see that $\dfdx{\theta}{s}=\pm\kappa$. From our diagram
we also see that
$$
\sin(\theta)=\dfdx{y}{s}\text{,  and }\cos(\theta)=\dfdx{x}{s}.
$$
\begin{embeddedproblem-enumerate}
         \begin{enumerate}[label={\bf{}(\alph*)}]{}
         \item Complete the line of reasoning we began in the
           previous paragraph to show that
           $\dx{y}=\pm\frac1\kappa\sin(\theta)\dx{\theta}$ and
           $\dx{x}=\pm\frac1\kappa\cos(\theta)\dx{\theta}$.
         \item Next show that
           \begin{align*}
             x\amp=\pm\frac1\kappa\sin(\theta) +x_0\\
             y\amp=\mp\frac1\kappa\cos(\theta) +y_0
           \end{align*}
           for some constants $x_0$ and $y_0$.
         \item And finally, show that the point $(x,y)$ lies on a
           circle of radius $\frac1\kappa$. What is the equation of the circle?
       \end{enumerate}

\end{embeddedproblem-enumerate}


%%% Local Variables: 
%%% outline-minor-mode: t
%%% eval:(flyspell-mode)
%%% TeX-master: "DiffCalc"
%%% End: 
